\section{Introduction}
\label{sec:intro}

Low-rank adaptation~\cite{hu2022lora} modifies pretrained language models
by injecting rank-$r$ update matrices $\Delta W = BA$ alongside frozen
weights. Multi-channel extensions~\cite{bielec2026weights,bielec2026bridge}
partition the rank dimension into $n$ channels governed by an
$n \times n$ learnable bridge matrix that controls inter-channel coupling.
The bridge adds $n^2$ parameters (36 at $n{=}6$) and introduces a degree
of freedom absent from standard LoRA: the \emph{topology} of the adapter.

Paper~3~\cite{bielec2026bridge} demonstrated that a closed-loop cybernetic
controller---the Steersman---drives 6-channel bridges toward 3-axis
block-diagonal structure aligned with the rhombic dodecahedron's (RD)
face-pair geometry. The result was consistent across six experiments,
42{,}500+ bridges, two model scales, and three initialization strategies,
with coupling ratios exceeding 22{,}000:1. Paper~3 framed this as
geometry \emph{discovery}: the Steersman finds structure latent in the
loss landscape.

This paper reframes the result. The bridge is not a discovery mechanism
but a \emph{programmable substrate}. The Steersman is not a geometer but
a \emph{topology compiler}. The pair specification---the declaration of
which channels are co-axial---is a program. Paper~3 ran one program (RD
face-pairs at $n{=}6$) and observed one output (three $2 \times 2$
blocks). The question this paper answers: does the compiler accept other
programs, and what happens when it receives an invalid one?

The answer is that the Steersman is general-purpose but selective. It
compiles any pair specification derived from polytope co-axial structure
and rejects specifications that lack genuine geometric coherence. With no
modification to the feedback mechanism, control laws, or training
hyperparameters---only the pair specification changes---the Steersman
programs three polytope geometries across two spatial dimensionalities:

\begin{itemize}
  \item \textbf{Octahedron} ($n{=}4$, 3D, 2 co-axial pairs):
    two $2 \times 2$ blocks, per-bridge mean coupling ratio
    \textbf{473{,}622:1}---the programme's strongest signal.
  \item \textbf{Rhombic dodecahedron} ($n{=}6$, 3D, 3 co-planar pairs):
    three $2 \times 2$ blocks, coupling ratio \textbf{70{,}404:1}.
    Paper~3's result, reproduced across seeds.
  \item \textbf{Tesseract} ($n{=}8$, 4D, 4 co-axial pairs):
    four $2 \times 2$ blocks, coupling ratio \textbf{41{,}564:1},
    Fiedler eigenvalue 0.000191, clean $4{+}4$ eigenvalue split.
\end{itemize}

\noindent Four negative controls define the compiler's rejection
boundary. Spectral-only training across $n \in \{3, 4, 8, 12\}$ reveals
a universal spectral attractor at Fiedler ${\approx}\,0.09$---the
bridge's natural connectivity equilibrium when feedback provides no
topological direction. Wrong-labels (WL-001: random 3-pair partition on
$n{=}6$) produces the predicted $3{+}3$ eigenvalue split but co/cross
ratio ${\sim}\,10^{-5}$:1---eigenvalue structure without directional
content. Resonance (R-001: number-theoretic pairs) produces chain-like
eigenvalues rather than block structure. An emanation architecture
(E-001: hierarchical master bridge with per-layer offsets) converges to
the spectral attractor, not to block-diagonal structure---hierarchy
provides connectivity without directionality.

These outcomes define four cleanly separated dynamical regimes with no
intermediate cases: block-diagonal ($\rho > 40{,}000$:1), spectral
attractor (Fiedler~${\approx}\,0.09$, $\rho \approx 1$:1),
hierarchical coherence ($\rho \approx 1$:1, Fiedler~${\approx}\,0.08$),
and connectivity collapse ($\rho \approx 10^{-5}$:1).

\medskip
\noindent We make the following contributions:

\begin{enumerate}
  \item \textbf{Multi-polytope topology programming.} The Steersman
    programs octahedral ($2{+}2$ at $n{=}4$), RD ($3{+}3$ at $n{=}6$),
    and tesseract ($4{+}4$ at $n{=}8$) geometry into bridge matrices
    using identical hyperparameters. Only the pair specification differs.

  \item \textbf{Spectral attractor universality.} Spectral-only
    training converges to Fiedler ${\approx}\,0.09$ across
    $n \in \{3, 4, 6, 8, 12\}$---a 19.4\% band across a $4\times$
    range of channel counts---establishing the bridge's unprogrammed
    equilibrium.

  \item \textbf{Geometric coherence requirement.} Wrong-labels control
    (random 3-pair partition on $n{=}6$) produces a clean $3{+}3$
    eigenvalue split with co/cross~${\sim}\,10^{-5}$:1, demonstrating
    that eigenvalue structure alone does not constitute topology
    programming. Directional content requires geometrically coherent
    pair specifications.

  \item \textbf{Number-theoretic negative result.} Prime-derived pairs
    (Sophie Germain, twin primes, residue classes) produce chain-like
    eigenvalue progressions, not block structure. Algebraic structure
    without geometric embedding is functionally equivalent to random
    noise ($<10\%$ trajectory divergence from wrong-labels across
    10{,}000 steps).

  \item \textbf{Emanation architecture.} A hierarchical master-plus-offsets
    bridge converges to the spectral attractor rather than block-diagonal
    structure, establishing that topology programming requires direct
    contrastive feedback to individual bridge matrices.

  \item \textbf{Reproducibility.} Tesseract contrastive ($n{=}8$)
    replicates with Pearson $r = 1.0000$ across 34 matching checkpoints
    from independent runs.

  \item \textbf{Four-regime taxonomy.} The complete experimental
    programme yields four dynamical regimes---block-diagonal, spectral
    attractor, wrong-partition, and collapse---cleanly separated with no
    intermediate outcomes.

  \item \textbf{Bridge initialization independence.} Three corpus-coupled
    initializations with maximally different sign patterns converge to
    statistically indistinguishable block-diagonal strength
    (50--52{,}000:1), confirming that topology is determined by the pair
    specification, not the initial conditions.
\end{enumerate}

The remainder of this paper is organized as follows.
\S\ref{sec:background} reviews Paper~3's findings and related work on
programmable inductive biases. \S\ref{sec:method} introduces the
generalized pair specification interface and the experimental design.
\S\ref{sec:three-geometries} presents the multi-polytope programming
results. \S\ref{sec:spectral-attractor} characterizes the spectral
attractor. \S\ref{sec:negative-controls} tests the compiler's rejection
boundary through wrong-labels and resonance controls.
\S\ref{sec:discussion} discusses implications, limitations, and the
four-regime taxonomy. \S\ref{sec:conclusion} concludes.
